\chapter{Liouville's Theorem and Motion in a Magnetic Field}

Liouville's theorem is one of the central facts of Hamiltonian mechanics. It
says that Hamiltonian evolution preserves volume in phase space. In quantum
mechanics the corresponding idea is unitarity: distinguishable states do not
quietly collapse into one another under the fundamental law of motion. In
both theories the underlying issue is the conservation of information.

We shall first make this geometric statement precise and prove it directly
from Hamilton's equations. We shall then use it to understand canonical area,
chaotic stretching, coarse-graining, and entropy. In the second half we turn
to a different problem, the motion of a charged particle in a magnetic field.
That example introduces a velocity-dependent force while preserving the
Lagrangian and Hamiltonian structures developed so far.

\section{A Flow That Does Not Lose Information}

A state of a classical system is a point in phase space. If we observe the
system at equally spaced instants, the law of motion draws an arrow from each
point to its successor. Determinism requires one outgoing arrow: the present
state must determine what happens next. Reversibility also requires one
incoming arrow: the present state must determine where it came from.

This already rules out two kinds of evolution. A trajectory cannot split into
two futures, and two distinct trajectories cannot merge into a single state.
Either event would destroy the one-to-one character of the microscopic law.
For a finite set of states, a reversible update rule is therefore a
permutation, decomposing the state space into disjoint cycles.

Hamiltonian mechanics makes a stronger statement. Fill a continuous phase
space uniformly with imaginary dust. Each dust grain follows the local
equations of motion. A small cloud of grains may stretch, bend, and shear, but
its phase-space volume remains unchanged. The flow is incompressible.

\begin{classroomqa}
\audiencequestion{Could a distribution of states converge or diverge even
when each individual trajectory is deterministic?\lecturetimestamp{00:05:05--00:06:15}}
\lecturerresponse{For a general deterministic rule that possibility must be
examined separately. Uniqueness of the future only forbids splitting.
Liouville's theorem is stronger: for a Hamiltonian system, an entire region
cannot be compressed into a smaller phase-space volume or expanded into a
larger one.}
\end{classroomqa}

There are several equivalent ways to picture the theorem. We may follow a
small volume surrounding each point; we may follow an arbitrarily shaped
region and compare its initial and final volumes; or we may begin with uniform
dust and demand that its density remain uniform. These are local, regional,
and density versions of the same incompressibility statement.

Volume preservation does not mean rigidity. A square can shear into a
parallelogram, or a nearly round cloud can become a long, thin filament. The
amount of phase space is fixed even though its shape and every ordinary
distance inside it may change.

\begin{figure}[tbp]
\centering
\includegraphics[width=0.78\textwidth]{lecture_07_figure_01.png}
\caption{A square and its sheared image at 00:13:16. Hamiltonian flow need
not preserve lengths or angles; the relevant invariant is phase-space area.}
\label{fig:l07-shear-board}
\end{figure}

\begin{classroomqa}
\audiencequestion{If the initial region is connected, can the flow tear it
apart or create a hole?\lecturetimestamp{00:10:08--00:11:11}}
\lecturerresponse{Not under a smooth Hamiltonian flow for which the solution
exists uniquely. The time-evolution map and its inverse are continuous, so
connectedness and topology are preserved. The region may become extremely
distorted, but it is not cut and reattached.}
\end{classroomqa}

\begin{classroomqa}
\audiencequestion{Can two points pass one another, and can the orientation of
a little square change?\lecturetimestamp{00:12:20--00:13:53}}
\lecturerresponse{Their projections onto one coordinate axis may pass. What
cannot happen is that two complete phase-space trajectories intersect at the
same time and state. A little square may rotate or shear into a
parallelogram; those changes are compatible with preserving its area.}
\end{classroomqa}

\begin{classroomqa}
\audiencequestion{Does a chaotic double pendulum contradict this picture?
Nearby initial conditions can separate rapidly.\lecturetimestamp{00:13:54--00:16:12}}
\lecturerresponse{No. Chaos concerns the growth of separations, not the
growth of exact phase-space volume. Expansion along unstable directions is
compensated by compression along other directions. A compact patch can
become a complicated filament while retaining precisely the same volume.}
\end{classroomqa}

For one canonical pair \((q,p)\), volume means area. For \(N\) pairs it means
the \(2N\)-dimensional volume element
\begin{equation}
 d\Gamma=\prod_{i=1}^{N}dq_i\,dp_i.
 \label{eq:l07-phase-volume-element}
\end{equation}
Two points are close only when their complete states are close: their
coordinates and their conjugate momenta must both be nearly equal.

\begin{classroomqa}
\audiencequestion{What exactly is the volume being preserved, and what does
``nearby'' mean?\lecturetimestamp{00:17:03--00:18:42}}
\lecturerresponse{For a single degree of freedom it is an area in the
\(q\)-\(p\) plane; with more degrees of freedom it is the product of all the
canonical coordinate and momentum intervals. Nearby states are close in this
full space, not merely close in ordinary position.}
\end{classroomqa}

Ordinary friction seems to violate the theorem. Many sliding states approach
the same visible state of rest, and many throws of a die can settle on the
same face. The apparent contraction arises because the description omits the
table, the air, and the internal degrees of freedom. Mechanical energy becomes
heat, sound, and microscopic motion. Information is transferred into those
unobserved variables. A reduced phase space may contract even though the
larger closed system remains Hamiltonian.

\begin{classroomqa}
\audiencequestion{When a die or a sliding block comes to rest, have many
states not genuinely merged?\lecturetimestamp{00:18:52--00:22:18}}
\lecturerresponse{They have merged only in a coarse description that records
the visible object and ignores its environment. Include the molecular state
of the table, the air, and the emitted sound, and the final microscopic states
remain distinct. The situation is analogous to energy: the block loses
kinetic energy while the table warms by a correspondingly tiny amount.}
\end{classroomqa}

\section{Incompressible Flow and the Divergence}

Before proving the theorem in phase space, consider an ordinary fluid of
points. In one dimension, let \(v_x(x)\) be the instantaneous velocity. A
small interval keeps a uniform density only if the flux entering one end
equals the flux leaving the other. Thus
\begin{equation}
 \frac{\partial v_x}{\partial x}=0.
 \label{eq:l07-one-dimensional-incompressibility}
\end{equation}
In one dimension an incompressible instantaneous flow is locally a rigid
translation.

In two dimensions take a small rectangle with sides \(\Delta x\) and
\(\Delta y\). The difference between outward fluxes through the right and
left faces is, to first order,
\begin{equation}
 \frac{\partial v_x}{\partial x}\,\Delta x\,\Delta y.
 \label{eq:l07-horizontal-flux}
\end{equation}
The corresponding difference through the upper and lower faces is
\begin{equation}
 \frac{\partial v_y}{\partial y}\,\Delta x\,\Delta y.
 \label{eq:l07-vertical-flux}
\end{equation}

\begin{figure}[tbp]
\centering
\begin{tikzpicture}[x=1.05cm,y=1.05cm,>=Stealth]
  \draw[thick] (0,0) rectangle (5,3);
  \draw[->,very thick,blue!65!black] (-1.15,0.75) -- (0,0.75)
    node[midway,above] {\(v_x(x)\)};
  \draw[->,very thick,blue!65!black] (-1.15,2.25) -- (0,2.25);
  \draw[->,very thick,blue!65!black] (5,0.75) -- (6.35,0.75)
    node[midway,above] {\(v_x(x+\Delta x)\)};
  \draw[->,very thick,blue!65!black] (5,2.25) -- (6.35,2.25);
  \draw[->,very thick,orange!75!black] (1.4,-1.05) -- (1.4,0)
    node[midway,left] {\(v_y(y)\)};
  \draw[->,very thick,orange!75!black] (3.6,-1.05) -- (3.6,0);
  \draw[->,very thick,orange!75!black] (1.4,3) -- (1.4,4.2)
    node[midway,left] {\(v_y(y+\Delta y)\)};
  \draw[->,very thick,orange!75!black] (3.6,3) -- (3.6,4.2);
  \draw[<->] (0,-0.38) -- (5,-0.38) node[midway,fill=white] {\(\Delta x\)};
  \draw[<->] (-0.38,0) -- (-0.38,3) node[midway,fill=white,rotate=90] {\(\Delta y\)};
\end{tikzpicture}
\caption{Instantaneous flux through a small rectangle. Incompressibility
requires the sum of the two first-order outward-flux differences to vanish.}
\label{fig:l07-divergence-reconstruction}
\end{figure}

After dividing by the area, incompressibility becomes
\begin{equation}
 \frac{\partial v_x}{\partial x}
 +\frac{\partial v_y}{\partial y}=0,
 \qquad\text{or}\qquad
 \bm{\nabla}\!\cdot\!\bm v=0.
 \label{eq:l07-two-dimensional-divergence}
\end{equation}
This does not require the velocity to be the same everywhere. An increase of
the \(x\)-velocity along \(x\) can be balanced by a decrease of the
\(y\)-velocity along \(y\). The rectangle may stretch in one direction and
contract in the other.

\begin{classroomqa}
\audiencequestion{A point entering through one side may turn and later leave
through another. Does the flux argument overlook that motion?
\lecturetimestamp{00:37:09--00:37:49}}
\lecturerresponse{The argument is instantaneous. Take a snapshot and ask for
the present rate of change of the number of points inside the rectangle. A
later turn belongs to a later snapshot. Summing the flux through all faces at
the same instant gives the local density change.}
\end{classroomqa}

The same reasoning works in any number of dimensions:
\begin{equation}
 \sum_{a=1}^{D}\frac{\partial v_a}{\partial x_a}=0.
 \label{eq:l07-general-divergence}
\end{equation}
For a nonuniform density \(\rho\), the general conservation law is the
continuity equation
\begin{equation}
 \frac{\partial\rho}{\partial t}
 +\bm{\nabla}\!\cdot(\rho\bm v)=0.
 \label{eq:l07-continuity-equation}
\end{equation}
Uniform dust reduces this to zero divergence.

\begin{classroomqa}
\audiencequestion{Does the argument depend on two dimensions or on a
uniform density?\lecturetimestamp{00:37:50--00:40:07}}
\lecturerresponse{The divergence formula extends to any number of
dimensions. Uniform density is the convenient case used here. For a varying
density one conserves the flux \(\rho\bm v\), as in
Eq.~\eqref{eq:l07-continuity-equation}.}
\end{classroomqa}

\section{Proof of Liouville's Theorem}

Return now to a phase space with canonical coordinates
\((q_1,\ldots,q_N,p_1,\ldots,p_N)\). Hamilton's equations are
\begin{equation}
 \dot q_i=\frac{\partial H}{\partial p_i},
 \qquad
 \dot p_i=-\frac{\partial H}{\partial q_i}.
 \label{eq:l07-hamilton-equations}
\end{equation}
They define a phase-space velocity field
\begin{equation}
 \bm u_H=(\dot q_1,\ldots,\dot q_N,
          \dot p_1,\ldots,\dot p_N).
 \label{eq:l07-hamiltonian-flow-field}
\end{equation}
Its divergence is
\begin{align}
 \bm{\nabla}_{q,p}\!\cdot\bm u_H
 &=\sum_{i=1}^{N}
 \left(
  \frac{\partial\dot q_i}{\partial q_i}
  +\frac{\partial\dot p_i}{\partial p_i}
 \right) \notag\\
 &=\sum_{i=1}^{N}
 \left(
  \frac{\partial^2 H}{\partial q_i\partial p_i}
  -\frac{\partial^2 H}{\partial p_i\partial q_i}
 \right)=0.
 \label{eq:l07-liouville-cancellation}
\end{align}
Each pair of mixed derivatives cancels. This local identity is the whole
proof.

\begin{figure}[tbp]
\centering
\includegraphics[width=0.82\textwidth]{lecture_07_figure_02.png}
\caption{The mixed-derivative cancellation at 00:43:55. The \(q_i\) and
\(p_i\) contributions enter with opposite signs.}
\label{fig:l07-liouville-board}
\end{figure}

\begin{theorem}[Liouville's theorem]
For a smooth Hamiltonian system, the time evolution preserves the canonical
phase-space volume element \(d\Gamma=\prod_i dq_i\,dp_i\).
\end{theorem}

\begin{proof}
Equation~\eqref{eq:l07-liouville-cancellation} says that the Hamiltonian
velocity field has zero divergence. The continuity equation then transports
an initially uniform density without changing it. Equivalently, the Jacobian
of the time-evolution map has unit determinant, so every transported region
keeps its phase-space volume.
\end{proof}

This is the promised strengthening of reversibility. A one-to-one flow need
not preserve volume in an arbitrary set of coordinates, but a Hamiltonian
flow in canonical coordinates does so because of the crossed derivatives and
the minus sign in Hamilton's equations.

\section{Canonical Area, Not Distance}

For one degree of freedom, Liouville volume is area in the \(q\)-\(p\)
plane. A simple change of coordinate shows why neither factor is separately
invariant. Take a free particle of unit mass,
\begin{equation}
 L=\frac{\dot x^2}{2},
 \qquad
 p_x=\frac{\partial L}{\partial\dot x}=\dot x.
 \label{eq:l07-free-particle}
\end{equation}
Introduce a rescaled coordinate
\begin{equation}
 y=\alpha x,
 \qquad
 x=\frac{y}{\alpha},
 \qquad
 \dot x=\frac{\dot y}{\alpha}.
 \label{eq:l07-coordinate-scaling}
\end{equation}
Then
\begin{equation}
 L=\frac{\dot y^2}{2\alpha^2},
 \qquad
 p_y=\frac{\partial L}{\partial\dot y}
 =\frac{\dot y}{\alpha^2}
 =\frac{p_x}{\alpha}.
 \label{eq:l07-scaled-momentum}
\end{equation}

\begin{figure}[tbp]
\centering
\includegraphics[width=0.80\textwidth]{lecture_07_figure_03.png}
\caption{The reciprocal scaling of coordinate and canonical momentum at
00:49:45: \(y=\alpha x\) while \(p_y=p_x/\alpha\).}
\label{fig:l07-scaling-board}
\end{figure}

The coordinate interval stretches while the conjugate momentum interval
shrinks:
\begin{equation}
 \Delta y\,\Delta p_y
 =(\alpha\Delta x)
 \left(\frac{\Delta p_x}{\alpha}\right)
 =\Delta x\,\Delta p_x.
 \label{eq:l07-area-invariance}
\end{equation}
This transformation is canonical. It changes the appearance of a cell but
not its symplectic area.

\begin{classroomqa}
\audiencequestion{Is \(\alpha\) a time-evolution parameter, and does every
coordinate transformation preserve area?\lecturetimestamp{00:54:24--00:55:14}}
\lecturerresponse{Here \(\alpha\) is only a fixed change of scale, such as a
change of units. The momentum must be transformed as the variable canonically
conjugate to the new coordinate. An arbitrary independent reshuffling of
\(q\) and \(p\) need not preserve area; a canonical transformation does.}
\end{classroomqa}

Hamiltonian time evolution is itself a canonical transformation. It can move
and deform a phase-space cell while preserving the canonical area element.
That statement is entirely compatible with energy conversion. A falling ball
trades potential energy for kinetic energy, so its position and momentum both
change while its total energy and the area of an ensemble are preserved.

\begin{classroomqa}
\audiencequestion{Does a falling ball violate conservation because its
momentum grows?\lecturetimestamp{00:53:27--00:54:23}}
\lecturerresponse{No. Potential energy decreases while kinetic energy
increases. Liouville's theorem concerns the area occupied by a collection of
nearby states, not the constancy of either coordinate or momentum for one
state.}
\end{classroomqa}

The same area already foreshadows quantum mechanics. Planck's constant has
the dimensions of position times momentum, and the uncertainty principle
sets a minimum scale for a quantum phase-space cell. The classical theorem
does not impose that minimum; it supplies the canonical geometry in which the
quantum statement will live.

\begin{classroomqa}
\audiencequestion{Why is Planck's constant associated with phase-space area
rather than with one axis?\lecturetimestamp{00:52:26--00:52:55}}
\lecturerresponse{Because its dimensions are action, \(q p\), and the
uncertainty relation constrains the product of conjugate uncertainties. A
rescaling can enlarge \(q\) while shrinking \(p\), but their canonical product
is the invariant quantity.}
\end{classroomqa}

\section{Chaos, Probability, and Coarse-Graining}

Liouville's theorem does not say that a phase-space patch remains simple. In
a chaotic system the patch may stretch exponentially, fold, and return across
the same macroscopic region in thinner and thinner filaments. Exact volume is
preserved throughout this process.

\begin{figure}[tbp]
\centering
\includegraphics[width=0.78\textwidth]{lecture_07_figure_04.png}
\caption{A compact patch and its filamented image at 01:03:40. The drawing
illustrates deformation under chaotic evolution, not an increase of exact
phase-space volume.}
\label{fig:l07-chaos-board}
\end{figure}

Phase-space area also supplies a natural invariant measure. Suppose a chaotic
trajectory explores an accessible component of an energy surface and define
probability by the fraction of a long observation time spent in a region.
When the motion is ergodic with respect to the Liouville measure, relative
occupation probabilities are proportional to invariant phase-space volume:
\begin{equation}
 \frac{P(\Omega_1)}{P(\Omega_2)}
 =\frac{\mu_L(\Omega_1)}{\mu_L(\Omega_2)}.
 \label{eq:l07-liouville-probability}
\end{equation}
Chaos by itself does not prove ergodicity, but it explains why this statistical
description can become useful even when the microscopic equations are fully
deterministic.

\begin{classroomqa}
\audiencequestion{Before quantum mechanics, what physical meaning does the
phase-space area have?\lecturetimestamp{00:57:50--00:59:43}}
\lecturerresponse{It supplies the invariant measure used in statistical
mechanics. Under the appropriate ergodic assumption, equal Liouville volumes
carry equal long-time statistical weight, so relative areas become relative
occupation probabilities.}
\end{classroomqa}

\begin{classroomqa}
\audiencequestion{Is this probability stochastic, or is chaotic mechanics
still deterministic?\lecturetimestamp{00:59:44--01:02:14}}
\lecturerresponse{The microscopic equations remain deterministic. The
statistical description enters because finite precision cannot track the
ever-finer distinctions among nearby initial conditions. Probability records
long-time occupation or incomplete knowledge, not an indeterministic force
inserted into the equations.}
\end{classroomqa}

To see the role of finite precision, partition phase space into cells of the
smallest size an apparatus can distinguish. These cells are fixed by the
measurement convention; they are not carried along and distorted with the
Hamiltonian flow. An initially compact patch may fit into a small set of cells.
After it has become finely filamented, it intersects many more fixed cells.

\begin{figure}[tbp]
\centering
\begin{tikzpicture}[x=0.78cm,y=0.78cm,>=Stealth]
  \begin{scope}
    \draw[step=0.75cm,gray!35,very thin] (0,0) grid (6,5);
    \fill[blue!35,draw=blue!70!black,thick]
      (2.0,1.6) .. controls (1.6,2.4) and (2.1,3.4) .. (3.1,3.5)
      .. controls (4.1,3.4) and (4.5,2.3) .. (4.0,1.5)
      .. controls (3.4,0.9) and (2.5,1.0) .. cycle;
    \node at (3,-0.55) {initial compact region};
  \end{scope}
  \begin{scope}[xshift=8cm]
    \draw[step=0.75cm,gray!35,very thin] (0,0) grid (6,5);
    \draw[blue!70!black,line width=5pt,rounded corners=7pt]
      (0.5,0.8) -- (1.1,3.9) -- (1.8,1.0) -- (2.5,4.3)
      -- (3.2,0.7) -- (3.9,4.0) -- (4.7,1.2) -- (5.5,4.2);
    \node at (3,-0.55) {same fine volume, larger cell cover};
  \end{scope}
  \draw[->,very thick] (6.35,2.5) -- (7.55,2.5) node[midway,above] {flow};
\end{tikzpicture}
\caption{Fine-grained volume and coarse-grained coverage. Stretching and
folding preserve exact volume, while finite-resolution cells can report a
larger occupied region.}
\label{fig:l07-coarse-graining-reconstruction}
\end{figure}

\begin{classroomqa}
\audiencequestion{Are the later coarse-graining cells the images of the
initial cells under time evolution?\lecturetimestamp{01:06:18--01:07:24}}
\lecturerresponse{No. The coarse cells represent what the apparatus can
distinguish at each time. If they were transported with the exact flow, they
would stretch into the same fine filaments and no coarse growth would appear.
The physically relevant comparison uses the same finite resolution before
and after evolution.}
\end{classroomqa}

The exact set still has its original volume. With infinitely precise
calculation and measurement, no expansion would be found. With fixed finite
resolution, however, the union of all cells intersecting the filaments can be
much larger. Looking with ``fuzzy eyes'' fills in separations too fine to
resolve.

\begin{classroomqa}
\audiencequestion{In what sense, then, can the phase-space area increase?
\lecturetimestamp{01:07:39--01:11:26}}
\lecturerresponse{Only the coarse-grained area increases. The fine-grained
set generated by Hamilton's equations retains exactly its original volume.
The larger value counts every resolution cell touched by the thin filaments,
including portions of those cells that the exact set does not occupy.}
\end{classroomqa}

This distinction points toward the second law of thermodynamics. Entropy
measures the logarithm of the detectable phase-space volume compatible with
our macroscopic knowledge. With a reference volume \(\Omega_0\), one may write
\begin{equation}
 S=k_B\log\!\left(\frac{\Omega_{\mathrm{cg}}}{\Omega_0}\right).
 \label{eq:l07-coarse-entropy}
\end{equation}
Microscopic information has not been destroyed. It has been hidden in
correlations and filaments finer than the available resolution.

\section{A Velocity-Dependent Fundamental Force}

We now turn to the motion of a particle of mass \(m\) and charge \(q\) in a
prescribed magnetic field. Its force is
\begin{equation}
 \boxed{\bm F=q\,\bm v\times\bm B.}
 \label{eq:l07-lorentz-force}
\end{equation}
Unlike the forces considered earlier, this one depends on velocity. Friction
also depends on velocity, but friction is an effective dissipative force in a
reduced description. The magnetic force is fundamental within the present
approximation and follows from an action principle.

The cross product has components
\begin{align}
 (\bm v\times\bm B)_x&=v_yB_z-v_zB_y,\notag\\
 (\bm v\times\bm B)_y&=v_zB_x-v_xB_z,\notag\\
 (\bm v\times\bm B)_z&=v_xB_y-v_yB_x.
 \label{eq:l07-cross-components}
\end{align}

\begin{figure}[tbp]
\centering
\includegraphics[width=0.82\textwidth]{lecture_07_figure_05.png}
\caption{Components of \(\bm v\times\bm B\) at 01:18:40. The upper line
isolates the \(z\)-component used in the derivation below.}
\label{fig:l07-cross-product-board}
\end{figure}

To construct a Lagrangian we introduce the vector potential \(\bm A(\bm x)\),
defined by
\begin{equation}
 \bm B=\bm\nabla\times\bm A.
 \label{eq:l07-vector-potential}
\end{equation}
In components,
\begin{align}
 B_x&=\frac{\partial A_z}{\partial y}
     -\frac{\partial A_y}{\partial z},\notag\\
 B_y&=\frac{\partial A_x}{\partial z}
     -\frac{\partial A_z}{\partial x},\notag\\
 B_z&=\frac{\partial A_y}{\partial x}
     -\frac{\partial A_x}{\partial y}.
 \label{eq:l07-curl-components}
\end{align}
Writing a field as a curl automatically gives \(\bm\nabla\cdot\bm B=0\), the
absence of magnetic monopole sources in the classical theory being used here.

The \(z\)-component of the magnetic force can now be expressed directly in
terms of \(\bm A\):
\begin{align}
 F_z
 &=q(v_xB_y-v_yB_x)\notag\\
 &=qv_x\left(
   \frac{\partial A_x}{\partial z}
  -\frac{\partial A_z}{\partial x}\right)
  +qv_y\left(
   \frac{\partial A_y}{\partial z}
  -\frac{\partial A_z}{\partial y}\right).
 \label{eq:l07-Fz-from-A}
\end{align}

\begin{figure}[tbp]
\centering
\includegraphics[width=0.82\textwidth]{lecture_07_figure_06.png}
\caption{The \(z\)-component of the magnetic force rewritten with
derivatives of the vector potential at 01:24:40.}
\label{fig:l07-force-from-A-board}
\end{figure}

\begin{classroomqa}
\audiencequestion{Why introduce \(\bm A\) when Newton's equation can be
written directly with \(\bm B\)?\lecturetimestamp{01:19:06--01:21:15}}
\lecturerresponse{The magnetic field is sufficient for the force law, but we
want an action, canonical momenta, and a Hamiltonian. The vector potential is
the quantity that couples locally to the velocity in the Lagrangian. Its curl
then makes the equations of motion depend on \(\bm B\).}
\end{classroomqa}

\section{The Magnetic Action and Canonical Momentum}

The free-particle action is supplemented by a line integral along the
particle's path:
\begin{equation}
 S=\int\frac12m v^2\,dt+q\int_{\gamma}\bm A\cdot d\bm x.
 \label{eq:l07-magnetic-action-line}
\end{equation}

\begin{figure}[tbp]
\centering
\includegraphics[width=0.82\textwidth]{lecture_07_figure_07.png}
\caption{The action as a kinetic term plus a line integral of the vector
potential at 01:28:50. The small path sketch emphasizes that the second term
depends on the route through space.}
\label{fig:l07-action-board}
\end{figure}

Since \(d\bm x=\dot{\bm x}\,dt\), the action takes the ordinary time-integral
form
\begin{equation}
 S=\int L\,dt,
 \qquad
 L=\frac12m\dot{\bm x}^{\,2}
   +q\,\bm A(\bm x)\cdot\dot{\bm x}.
 \label{eq:l07-magnetic-lagrangian}
\end{equation}
We use units in which the coupling is written \(q\bm A\cdot\bm v\). Other
electromagnetic unit conventions may place a factor of \(1/c\) in this term.

\begin{classroomqa}
\audiencequestion{How can the two terms in the action have the same units?
\lecturetimestamp{01:29:51--01:31:17}}
\lecturerresponse{The action has units of energy times time. In the chosen
convention \(q\bm A\) has units of momentum, so
\(q\bm A\cdot d\bm x\) has units of action. A different convention may move a
factor of \(c\) into the definition of the coupling, without changing the
mechanics.}
\end{classroomqa}

The canonical momentum is no longer simply the mechanical momentum:
\begin{equation}
 p_i=\frac{\partial L}{\partial\dot x_i}
 =m\dot x_i+qA_i.
 \label{eq:l07-canonical-momentum}
\end{equation}
Thus
\begin{equation}
 \bm p_{\mathrm{mech}}=m\bm v=\bm p-q\bm A.
 \label{eq:l07-mechanical-momentum}
\end{equation}

\begin{figure}[tbp]
\centering
\includegraphics[width=0.78\textwidth]{lecture_07_figure_08.png}
\caption{Canonical momentum in the \(x\)-direction at 01:35:20:
\(p_x=m\dot x+qA_x\).}
\label{fig:l07-canonical-momentum-board}
\end{figure}

This distinction is essential. Hamilton's equations use the canonical
momentum \(\bm p\); kinetic energy and the Lorentz force are most naturally
expressed through the mechanical combination \(\bm p-q\bm A\).

\section{Recovering the Lorentz Force}

We must verify that the proposed Lagrangian gives the correct equation of
motion. Work out the \(z\)-component explicitly. The Euler--Lagrange equation
is
\begin{equation}
 \frac{d}{dt}\frac{\partial L}{\partial\dot z}
 =\frac{\partial L}{\partial z}.
 \label{eq:l07-z-euler-lagrange}
\end{equation}
The left-hand side is
\begin{equation}
 \frac{d}{dt}(m\dot z+qA_z)
 =m\ddot z+q\frac{dA_z}{dt}.
 \label{eq:l07-z-left-side}
\end{equation}
The right-hand side receives contributions from all three components of the
vector potential:
\begin{equation}
 \frac{\partial L}{\partial z}
 =q\dot x\frac{\partial A_x}{\partial z}
 +q\dot y\frac{\partial A_y}{\partial z}
 +q\dot z\frac{\partial A_z}{\partial z}.
 \label{eq:l07-z-right-side}
\end{equation}

The field is assumed static, meaning it has no explicit time dependence. It
may still vary from point to point. Because the particle moves through that
spatially varying field,
\begin{equation}
 \frac{dA_z}{dt}
 =\dot x\frac{\partial A_z}{\partial x}
 +\dot y\frac{\partial A_z}{\partial y}
 +\dot z\frac{\partial A_z}{\partial z}.
 \label{eq:l07-Az-chain-rule}
\end{equation}

\begin{classroomqa}
\audiencequestion{If the magnetic field is static, why is
\(dA_z/dt\) not zero?\lecturetimestamp{01:39:13--01:39:53}}
\lecturerresponse{Static means no explicit dependence on time, not uniform
in space. The particle samples different values of \(\bm A(\bm x)\) as it
moves, so the value of \(A_z\) along the trajectory changes by the spatial
chain rule.}
\end{classroomqa}

\begin{figure}[tbp]
\centering
\includegraphics[width=0.84\textwidth]{lecture_07_figure_09.png}
\caption{The two sides of the \(z\)-component Euler--Lagrange equation at
01:41:50. Terms proportional to
\(\dot z\,\partial A_z/\partial z\) occur on both sides and cancel.}
\label{fig:l07-euler-lagrange-board}
\end{figure}

Substitute Eq.~\eqref{eq:l07-Az-chain-rule} into the left-hand side of
Eq.~\eqref{eq:l07-z-euler-lagrange}. The terms
\(q\dot z\,\partial A_z/\partial z\) cancel, leaving
\begin{align}
 m\ddot z
 &=q\dot x\left(
   \frac{\partial A_x}{\partial z}
  -\frac{\partial A_z}{\partial x}\right)
  +q\dot y\left(
   \frac{\partial A_y}{\partial z}
  -\frac{\partial A_z}{\partial y}\right)\notag\\
 &=q(v_xB_y-v_yB_x)
 =q(\bm v\times\bm B)_z.
 \label{eq:l07-lorentz-z-derived}
\end{align}
The other components follow by cyclic permutation, and therefore
\begin{equation}
 m\ddot{\bm x}=q\,\dot{\bm x}\times\bm B.
 \label{eq:l07-full-lorentz-derived}
\end{equation}

\begin{figure}[tbp]
\centering
\includegraphics[width=0.84\textwidth]{lecture_07_figure_10.png}
\caption{The surviving derivative combinations at 01:44:20, identified as
\(B_y\) and \(-B_x\), complete the \(z\)-component of the Lorentz force.}
\label{fig:l07-lorentz-recovery-board}
\end{figure}

\begin{classroomqa}
\audiencequestion{Is the vector potential uniquely determined by the
magnetic field?\lecturetimestamp{01:44:34--01:45:10}}
\lecturerresponse{No. For a static gauge transformation
\(\bm A\mapsto\bm A+\bm\nabla\chi\), the curl is unchanged. The Lagrangian
changes by the total derivative \(q\,d\chi/dt\), which changes the action only
by endpoint terms and leaves the equations of motion unchanged.}
\end{classroomqa}

The contrast with friction is now sharp. A frictional force has a component
opposite the velocity and changes the speed. A magnetic force is perpendicular
to both \(\bm v\) and \(\bm B\):
\begin{equation}
 \bm F\cdot\bm v
 =q(\bm v\times\bm B)\cdot\bm v=0.
 \label{eq:l07-magnetic-power-zero}
\end{equation}
It changes the direction of motion without changing the speed.

\begin{classroomqa}
\audiencequestion{Why is a velocity-dependent magnetic force so different
from friction?\lecturetimestamp{01:45:15--01:46:28}}
\lecturerresponse{The decisive issue is direction. Friction acts along the
velocity axis and therefore changes kinetic energy. The magnetic force is
perpendicular to the velocity, so its instantaneous power vanishes. It bends
the trajectory but does not speed the particle up or slow it down.}
\end{classroomqa}

\section{Energy and the Hamiltonian}

Assume that \(\bm A\) has no explicit time dependence. The Hamiltonian is
\begin{equation}
 H=\sum_i p_i\dot x_i-L.
 \label{eq:l07-H-definition}
\end{equation}
First express it in terms of velocities:
\begin{align}
 H
 &=\sum_i(m\dot x_i+qA_i)\dot x_i
 -\left(
   \frac12m\sum_i\dot x_i^2
   +q\sum_iA_i\dot x_i
  \right)\notag\\
 &=\frac12m\sum_i\dot x_i^2
 =\frac12mv^2.
 \label{eq:l07-H-velocity}
\end{align}
The two copies of the magnetic coupling cancel.

\begin{figure}[tbp]
\centering
\includegraphics[width=0.84\textwidth]{lecture_07_figure_11.png}
\caption{The \(x\)-component of the Legendre transform at 01:54:20. The
\(qA_x\dot x\) terms cancel, leaving \(m\dot x^2/2\).}
\label{fig:l07-energy-cancellation-board}
\end{figure}

This agrees with the force argument: a static magnetic field does no work.
In a uniform field perpendicular to the velocity, the particle moves in a
circle at constant speed. Its direction continually changes, so it has an
acceleration, but its kinetic energy remains constant.

To use Hamilton's equations we must eliminate the velocities. From
Eq.~\eqref{eq:l07-canonical-momentum},
\begin{equation}
 \dot x_i=\frac{p_i-qA_i}{m}.
 \label{eq:l07-velocity-from-canonical-momentum}
\end{equation}

\begin{figure}[tbp]
\centering
\includegraphics[width=0.82\textwidth]{lecture_07_figure_12.png}
\caption{The Hamiltonian definition and the inversion
\(\dot x=(p_x-qA_x)/m\) at 01:57:00.}
\label{fig:l07-H-inversion-board}
\end{figure}

The Hamiltonian in canonical variables is therefore
\begin{equation}
 \boxed{
 H(\bm x,\bm p)
 =\frac{1}{2m}\bigl(\bm p-q\bm A(\bm x)\bigr)^2.}
 \label{eq:l07-magnetic-Hamiltonian}
\end{equation}
It looks more complicated than the kinetic energy because the canonical
momentum contains the vector potential. The combination
\(\bm p-q\bm A=m\bm v\) is the mechanical momentum.

\begin{classroomqa}
\audiencequestion{An accelerated charge radiates. How can its kinetic energy
remain constant?\lecturetimestamp{01:57:29--01:58:29}}
\lecturerresponse{Radiation is being neglected in this
particle-in-an-external-field model. A complete treatment must include the
electromagnetic field as a dynamical Hamiltonian system carrying its own
energy. Radiation reaction is a higher-order correction to the approximation
used here.}
\end{classroomqa}

As a final check, Hamilton's equations give
\begin{equation}
 \dot x_i=\frac{\partial H}{\partial p_i}
 =\frac{p_i-qA_i}{m}=v_i,
 \label{eq:l07-first-H-check}
\end{equation}
and
\begin{equation}
 \dot p_i=-\frac{\partial H}{\partial x_i}
 =qv_j\frac{\partial A_j}{\partial x_i}.
 \label{eq:l07-second-H-check}
\end{equation}
Differentiating the mechanical momentum gives
\begin{align}
 m\dot v_i
 &=\dot p_i-q\frac{dA_i}{dt}\notag\\
 &=qv_j\left(
   \frac{\partial A_j}{\partial x_i}
  -\frac{\partial A_i}{\partial x_j}
  \right)
 =q(\bm v\times\bm B)_i.
 \label{eq:l07-Hamiltonian-lorentz-check}
\end{align}
Thus Newtonian, Lagrangian, and Hamiltonian descriptions all produce the same
magnetic motion. The vector potential changes the canonical bookkeeping, but
the observable force depends only on its curl.
